% Удиртгал
% \phantomchapter{Удиртгал} % Зарим нэг зөвлөмж
\phantomsection
\addchaptertocentry{Удиртгал}
\label{Introduction} % Энэ бүлэг рүү ишлэл хийх бол \ref{Introduction} командыг ашигла
%-------------------------------------------------------------------------------
%   INTRODUCTION
%-------------------------------------------------------------------------------
% Зорилго зорилт, өмнөх бие даалт харах
    \section*{Удиртгал}
    Орчин үеийн ухаалаг хотын хөгжлийн нэг чухал чиглэл нь иргэдийн оролцоонд тулгуурласан мэдээллийн системүүд юм. Ялангуяа дугуйн хэрэглээ нэмэгдэж буй өнөөгийн нөхцөлд дугуйн замын аюулгүй байдал, хүртээмж болон замын нөхцлийн мэдээллийг бодит цагт цуглуулах, боловсруулах шаардлага улам бүр нэмэгдэж байна. Гэвч одоогийн байдлаар дугуйн замын мэдээлэл нь төвлөрсөн бус, тогтмол шинэчлэгддэггүй, хэрэглэгчдэд бүрэн хүрч чаддаггүй асуудал түгээмэл байна.
    Иймээс олон хэрэглэгчийн өгөгдөлд тулгуурлан замын нөхцлийг тодорхойлох, аюулын цэгүүдийг илрүүлэх, мөн ухаалаг маршрутын аюулгүй байдлыг визуалчлах систем хөгжүүлэх нь практик ач холбогдол өндөртэй гэж үзэж байна. Энэхүү систем нь хэрэглэгчдийн оруулсан мэдээллийг нэгтгэн боловсруулж, дугуйн замын нөхцлийг илүү бодитой, динамик байдлаар харуулах боломжийг бүрдүүлнэ.

    \subsection*{Зорилго}
    Энэхүү төслийн ажлын зорилго нь Улаанбаатар хотын дугуйн замын мэдээллийг олон нийтийн оролцоонд тулгуурлан цуглуулах, боловсруулах, мөн замын нөхцөл болон аюулын мэдээллийг газрын зураг дээр визуалчлах боломжтой crowdsourcing веб системийг хөгжүүлэхэд оршино.

    \subsection*{Зорилт}
    Уг зорилгод хүрэхийн тулд дараах зорилтуудыг дэвшүүлж байна. Үүнд:
    \begin{enumerate}[nosep]
        \item Crowdsourcing болон GIS суурьтай ижил төстэй системүүдийг судлах;
        \item Дугуйн замын нөхцөл болон POI системийн өгөгдлийн загварыг тодорхойлох;
        \item Хэрэглэгч газрын зураг дээр замын сегмент тэмдэглэх (ногоон/шар/улаан) боломжтой систем хөгжүүлэх;
        \item Аюулын болон сонирхолтой цэг (POI) нэмэх, санал өгөх, удирдах функц хэрэгжүүлэх;
        \item Олон хэрэглэгчийн саналыг spatial hash -ээр нэгтгэн (crowd aggregation) тухайн замын давамгай нөхцлийг тодорхойлох алгоритм боловсруулах;
        \item OSRM маршрут дээр crowd өнгөт давхарга болон аюулын POI -г хослуулсан Smart Route визуалчлал хэрэгжүүлэх;
        \item GPS ашиглан явсан замыг session -д бичиж GPX файл болгож татах боломжийг олгох;
        \item Админд POI -г батлах/татгалзах болон системийн статистик хянах хэрэгсэл бүрдүүлэх;
        \item Хөгжүүлсэн системийн үр дүн болон ажиллагаанд дүгнэлт хийх.
    \end{enumerate}

    \subsection*{Судлах үндэслэл}
        Crowdsourcing болон газрын зурагт суурилсан ухаалаг системүүд нь олон хэрэглэгчийн өгөгдлийг бодит цагт цуглуулж, боловсруулан шийдвэр гаргалтад ашиглах боломжийг олгодог тул орчин үеийн ухаалаг хотын хөгжлийн чухал чиглэлд тооцогдож байна. Ялангуяа дугуйн замын мэдээлэл нь ихэвчлэн төвлөрсөн эх сурвалжаас хамааралтай, шинэчлэл нь удаан байдаг тул хэрэглэгчдийн оролцоонд тулгуурласан систем хөгжүүлэх шаардлага бий болж байна.
        Энэхүү судалгааны үндэслэл нь дараах давуу талуудтай холбоотой юм:
        \begin{enumerate}[nosep]
            \item Олон хэрэглэгчийн өгөгдлөөр илүү бодит мэдээлэл бий болгоно
            \par Хэрэглэгч бүрийн оруулсан замын нөхцөл болон POI мэдээллийг нэгтгэснээр ганц эх сурвалжаас илүү бодит, өргөн хүрээтэй мэдээллийн сан бүрдэнэ.
            \item Замын аюулгүй байдлын үнэлгээг сайжруулна
            \par Crowd aggregation аргаар олон хэрэглэгчийн саналд тулгуурлан замын сегментүүдийг ногоон (дугуйн замтай, аюулгүй), шар (боломжтой), улаан (аюултай буюу тохиромжгүй) гэсэн ангиллаар тодорхойлох боломжтой бөгөөд энэ нь илүү найдвартай үнэлгээ гаргана.
            \item Маршрутын аюулгүй байдлын мэдээллийг харуулна
            \par OSRM маршрутын дээр POI болон сегментэд суурилсан өгөгдлийг давхаргалснаар хэрэглэгч аюулгүй замаа сонгоход дэмжлэг өгнө.
            \item Хотын дэд бүтцийн төлөвлөлтөд дэмжлэг үзүүлнэ
            \par Цугларсан өгөгдөл нь дугуйн зам дутмаг эсвэл аюултай хэсгүүдийг тодорхойлох боломж олгож, хот төлөвлөлт болон замын дэд бүтцийн сайжруулалтад ашиглагдана.
        \end{enumerate}

    \subsection*{Судлагдсан байдал}
        Crowd-sourcing болон газрын зурагт суурилсан ухаалаг системүүд нь олон хэрэглэгчийн өгөгдлийг цуглуулж, боловсруулан нэгтгэх замаар илүү бодит мэдээллийн сан бүрдүүлэх боломжийг олгодог. Ийм төрлийн системүүд нь хэрэглэгчдийн оруулсан мэдээлэлд тулгуурлан газрын зургийн мэдээллийг байнга шинэчилж, илүү динамик байдлаар ашиглах нөхцөл бүрдүүлдэг.
        Жишээлбэл, хэрэглэгчид өөрсдийн явсан замын хэсгүүдэд нөхцлийн тэмдэглэгээ хийх, аюултай цэгүүдийг газрын зураг дээр нэмэх, мөн бусад хэрэглэгчдийн оруулсан мэдээллийг харах боломжтой байдаг. Энэ нь төвлөрсөн нэг эх сурвалжаас хамааралгүйгээр мэдээллийг олон эх үүсвэрээс цуглуулах давуу талтай.
        Мөн олон хэрэглэгчийн өгөгдлийг нэгтгэх (crowd aggregation) аргыг ашигласнаар тухайн замын хэсгийн хамгийн их давтамжтай саналд тулгуурлан бодит нөхцлийг тодорхойлох боломжтой болдог. Ийм аргыг ашигласнаар замын аюулгүй байдлын үнэлгээ илүү найдвартай, бодитой болдог.
        Орчин үеийн crowd-sourcing суурьтай системүүдийн нэг жишээ нь Strava зэрэг хэрэглэгчийн хөдөлгөөний өгөгдөлд тулгуурлан маршрут болон замын мэдээлэл үүсгэдэг платформууд юм. Мөн OpenStreetMap нь олон нийтийн оролцоонд тулгуурласан газрын зургийн нээлттэй платформ бөгөөд OSRM (Open Source Routing Machine) нь түүний өгөгдлийг ашиглан маршрут тооцоолохад ашиглагддаг.

    \subsection*{Шинэлэг тал}
    \textbf{Crowdsourcing дугуйн замын веб системийн шинэлэг тал:}
        \begin{enumerate}[itemsep=0mm]
            \item Олон хэрэглэгчийн оролцоонд тулгуурлан замын мэдээллийг бодит цагт цуглуулах боломжтой;
            \item Замын сегмент бүрийг нөхцлөөр нь ангилан тэмдэглэх (ногоон, шар, улаан) болон 1-6 түвшний дэд бүтцийн зэрэглэлээр тэмдэглэх боломж олгосон;
            \item Аюулын болон сонирхолтой цэгүүд (POI)-ийг 6 төрлөөр (Danger, No Bike Lane, Road Damage, Parking Problem, Bike Repair, Bike Parking) газрын зураг дээр тэмдэглэж, upvote/downvote санал өгөх, админ батлах боломжтой;
            \item \textbf{Crowd Aggregation} -- Олон хэрэглэгчийн саналыг spatial hash (SHA-256, 110м нарийвчлал) -ээр нэгтгэн тухайн замын давамгай нөхцлийг $\arg\max$ үйлдлээр тодорхойлох алгоритм ашигласан;
            \item \textbf{Interactive map visualization} -- Leaflet + OpenStreetMap дээр замын нөхцөл, POI мэдээллийг ойлгомжтой байдлаар визуалчилсан;
            \item \textbf{User-generated data} -- Хэрэглэгч өөрөө өгөгдөл үүсгэж, системийн мэдээллийн санг өргөжүүлэх боломжтой;
            \item \textbf{Smart Route annotation} -- OSRM cycling маршрутын дээр crowd aggregation -ээс гаргасан өнгөт давхарга болон ойролцоох (\textasciitilde 55м) аюулын POI -г hazard layer хэлбэрээр давхаргалсан;
            \item \textbf{GPS \& GPX support} -- GPS ашиглан явсан замыг browser session -д бодит цагт бичих, нийт явсан километрийг профайлд бүртгэх, GPX 1.1 стандартын дагуу .gpx файл болгож татах боломжтой.
        \end{enumerate}

    \subsection*{Технологийн ач холбогдол}
         Дээрх шинэлэг талуудаас энэхүү систем нь дараах ач холбогдлуудыг авчирна.
                 \begin{enumerate}[itemsep=0mm]
            \item Дугуйн замын бодит нөхцлийн мэдээллийг хурдан бөгөөд хүртээмжтэй болгох
            \item Хэрэглэгчдийн оролцоонд суурилсан нээлттэй мэдээллийн сан бүрдүүлэх
            \item Маршрут дээрх аюулгүй байдлын мэдээллийг визуалчилж сонголт хийхэд дэмжлэг үзүүлэх
            \item Хотын дэд бүтцийн төлөвлөлтөд ашиглагдах бодит өгөгдөл бий болгох
            \item Server-Side Rendering (Django Templates + Vanilla JavaScript) хандлагаар хөнгөн, build алхамгүй веб ажиллагааг хангах
        \end{enumerate}

    \subsection*{Хамрах хүрээ}
    Энэхүү crowd-sourcing дугуйн замын веб систем нь дугуй унадаг иргэд, сонирхогчид болон хотын замын нөхцлийн мэдээлэл сонирхдог хэрэглэгч бүр ашиглах боломжтой. Хэрэглэгчид газрын зураг дээр замын нөхцөл тэмдэглэх, аюулын болон сонирхолтой цэг (POI) нэмэх, бусад хэрэглэгчдийн оруулсан мэдээллийг харах, A→B аюулгүй маршрутыг визуалчлан харах боломжтой. Систем нь зочин (бүртгэлгүй), хэрэглэгч (бүртгэлтэй) болон админ гэсэн 3 төрлийн ролетой бөгөөд RBAC -аар эрхийг хязгаарлана.
    Мөн уг систем нь хот төлөвлөлт, замын дэд бүтцийн судалгаа хийдэг байгууллагууд болон судлаачдад бодит өгөгдөлд тулгуурласан мэдээллийн эх сурвалж болох боломжтой. Веб технологид суурилсан тул интернэт холболттой ямар ч төхөөрөмжөөс (гар утас, ширээний компьютер) хандаж Bootstrap 5 -ийн responsive UI -аар ашиглах боломжтой юм.
